% concatenated from generic_equivalence_arxiv.tex
\documentclass{amsart}
\usepackage{amssymb,
enumitem,
mathrsfs,
mathtools,
tikz,
upgreek,
verbatim,
hyphenat,
url
}
\usepackage[T1]{fontenc}
\usepackage[shadow]{todonotes}
\usepackage{tikz-cd}
\renewcommand{\baselinestretch}{1.2}
\usepackage{newpxtext,newpxmath}
\usepackage{subfig}
\usepackage{bbold}
\usetikzlibrary{calc}

\usetikzlibrary{shapes.misc, positioning}


\usepackage[colorlinks=true,linkcolor={black},citecolor={blue}]{hyperref}%

% other fonts

\newcommand{\pre}[2]{{}^{#1} #2}
\newcommand{\set}[1]{\{ #1 \}}
\newcommand{\setm}[2]{\{ #1 : #2 \}}
\newcommand{\On}{\mathrm{On}}
\newcommand{\markdef}[1]{\textbf{#1}}

% operators
\newcommand{\id}{\operatorname{id}}
\newcommand{\pow}{\mathscr{P}}
\newcommand{\Mod}{\operatorname{Mod}}
\newcommand{\lh}{\operatorname{lh}}
\newcommand{\pred}{\operatorname{pred}}
\newcommand{\dom}{\operatorname{dom}}
\newcommand{\cod}{\operatorname{cod}}
\newcommand{\Gen}{\operatorname{Gen}}
\newcommand{\cat}{^\frown}

\def\colim{\qopname\relax m{colim}}


\newcommand{\cf}{{\operatorname{cf}}}
\newcommand{\inj}{{\operatorname{inj}}}
\newcommand{\rk}{{\operatorname{rk}}}
\newcommand{\MS}{\mathfrak{M}}
\newcommand{\D}[2]{#1_{#2\uparrow}}
\newcommand{\lev}{\operatorname{lev}}

% axioms
\newcommand{\ZF}{{\sf ZF}}
\newcommand{\AC}{{\sf AC}}

% others
\newcommand{\I}[2]{I_{#1}^{\vec{#2}}}
\newcommand{\II}[1]{I^{\vec{#1}}}
\newcommand{\T}[2]{\mathbb{T}_{#1}^{#2}}
\newcommand{\Isomm}{\cong}
\newcommand{\RCT}{\mathrm{RCT}}
\newcommand{\embeds}{\sqsubseteq}
\newcommand{\analytic}{\boldsymbol{\Sigma_1^1}}
\newcommand{\Subg}[1]{\mathrm{Subg}(#1)}
\newcommand{\Isom}[1]{\mathrm{Isom}(#1)}
\renewcommand{\wr}{\mathbin{\mathrm{wr}}}
\newcommand{\Fix}{\mathrm{Fix}}
\newcommand{\LO}{\mathrm{LO}}
\newcommand{\CLO}{\mathrm{CLO}}
\newcommand{\CO}{\mathrm{CO}}
\newcommand{\Diff}{\mathrm{Diff}}


\DeclareMathOperator{\rank}{rank}
\DeclareMathOperator{\proj}{proj}
\DeclareMathOperator{\GL}{GL}
\DeclareMathOperator{\SL}{SL}
\DeclareMathOperator{\OO}{O}
\DeclareMathOperator{\SO}{SO}
\DeclareMathOperator{\Spin}{Spin}
\DeclareMathOperator{\PSO}{PSO}
\DeclareMathOperator{\PGL}{PGL}
\DeclareMathOperator{\Aut}{Aut}
\DeclareMathOperator{\BS}{BS}
\DeclareMathOperator{\Fr}{Fr}
\DeclareMathOperator{\Rad}{Rad}
\DeclareMathOperator{\spn}{span}
\DeclareMathOperator{\Ar}{Ar}
\DeclareMathOperator{\Homeo}{Homeo}
\DeclareMathOperator{\Orb}{Orb}
\DeclareMathOperator{\Stab}{Stab}
\DeclareMathOperator{\Add}{Add}
\DeclareMathOperator{\graph}{graph}






\newcommand\N{\mathbb{N}}
\newcommand\Z{\mathbb{Z}}
\newcommand\Q{\mathbb{Q}}
\newcommand\R{\mathbb{R}}
\renewcommand\S{\mathbb{S}}









% lists
\newenvironment{enumerate-(a)}{\begin{enumerate}[label={\upshape (\alph*)}, leftmargin=2pc]}{\end{enumerate}}
\newenvironment{enumerate-(a)-r}{\begin{enumerate}[label={\upshape (\alph*)}, leftmargin=2pc,resume]}{\end{enumerate}}
\newenvironment{enumerate-(a)-5}{\begin{enumerate}[label={\upshape (\alph*)}, leftmargin=2pc,start=5]}{\end{enumerate}}
\newenvironment{enumerate-(A)}{\begin{enumerate}[label={\upshape (\Alph*)}, leftmargin=2pc]}{\end{enumerate}}
\newenvironment{enumerate-(A)-r}{\begin{enumerate}[label={\upshape (\Alph*)}, leftmargin=2pc,resume]}{\end{enumerate}}
\newenvironment{enumerate-(i)}{\begin{enumerate}[label={\upshape (\roman*)}, leftmargin=2pc]}{\end{enumerate}}
\newenvironment{enumerate-(i)-r}{\begin{enumerate}[label={\upshape (\roman*)}, leftmargin=2pc,resume]}{\end{enumerate}}
\newenvironment{enumerate-(I)}{\begin{enumerate}[label={\upshape (\Roman*)}, leftmargin=2pc]}{\end{enumerate}}
\newenvironment{enumerate-(I)-r}{\begin{enumerate}[label={\upshape (\Roman*)}, leftmargin=2pc,resume]}{\end{enumerate}}
\newenvironment{enumerate-(1)}{\begin{enumerate}[label={\upshape (\arabic*)}, leftmargin=2pc]}{\end{enumerate}}
\newenvironment{enumerate-(1)-r}{\begin{enumerate}[label={\upshape (\arabic*)}, leftmargin=2pc,resume]}{\end{enumerate}}
\newenvironment{itemizenew}{\begin{itemize}[leftmargin=2pc]}{\end{itemize}}

% theorems
\newtheorem{theorem}{Theorem}[section]
\newtheorem{lemma}[theorem]{Lemma}
\newtheorem{corollary}[theorem]{Corollary}
\newtheorem{proposition}[theorem]{Proposition}
\newtheorem{question}[theorem]{Question}
\newtheorem{problem}[theorem]{Problem}
\newtheorem{fact}[theorem]{Fact}
\newtheorem{claim}{Claim}[theorem]
\theoremstyle{definition}
\newtheorem{definition}[theorem]{Definition}
\newtheorem{notation}[theorem]{Notation}
\newtheorem{coding}{Coding}
\newtheorem{example}[theorem]{Example}
\newtheorem{conjecture}[theorem]{Conjecture}
\theoremstyle{remark}
\newtheorem{remark}[theorem]{Remark}



\newtheorem{innercustomgeneric}{\customgenericname}
\providecommand{\customgenericname}{}
\newcommand{\newcustomtheorem}[2]{%
  \newenvironment{#1}[1]
  {%
   \renewcommand\customgenericname{#2}%
   \renewcommand\theinnercustomgeneric{##1}%
   \innercustomgeneric
  }
  {\endinnercustomgeneric}
}

\newcustomtheorem{customthm}{Theorem}
\newcustomtheorem{customlemma}{Lemma}




\title{Forcing, genericity, and CBERS}



\author[F.~Calderoni]{Filippo Calderoni}

\address{Department of Mathematics, Rutgers University, Hill Center for the Mathematical Sciences, 110 Frelinghuysen Rd., Piscataway, NJ 08854-8019} \email{filippo.calderoni@rutgers.edu}

\author[D.~Sinapova]{Dima Sinapova}
\address{Department of Mathematics, Rutgers University, Hill Center for the Mathematical Sciences, 110 Frelinghuysen Rd., Piscataway, NJ 08854-8019} \email{dima.sinapova@rutgers.edu}


\subjclass[2020]{03E15, 03E40}

\thanks{The research of the first author was supported by the NSF grant DMS -- 2348819. The research of the second author was supported by the NSF grant DMS -- 2246781.}



\begin{document}

\maketitle

\begin{abstract}
    In this paper we continue the study of equivalence of generics filters started by Smythe in \cite{Smy22}. We fully characterize those forcing posets for which the corresponding equivalence of generics is smooth using the purely topological property of condensation. Next we leverage our characterization to show that there are non-homogeneous forcing for which equivalence of generics is not smooth. Then we prove hyperfiniteness in the case of Prikry forcing and present some additional results addressing the problem of whether generic equivalence for Cohen forcing is hyperfinite.
\end{abstract}

\section{Introduction}

Let \(M\) be a countable transitive model of a sufficient fragment of \(\textsf{ZFC}\), and fix a partial order
\(\mathbb{P}\) in \(M\).
We can define the Polish space $\Gen(M, \mathbb{P})$ of $M$-generic filters for $\mathbb{P}$.  This space has a basis of clopen sets of the form $N_p=\{G\in \Gen(M, \mathbb{P})\mid p\in G\}$ for all $p\in \mathbb{P}\cap M$.
We say that two generic filters \(G,H\in \Gen(M, \mathbb{P})\)
are \markdef{equivalent} (in symbols, \(G \equiv_{M}^\mathbb{P} H\)) if and only if they produce the same generic extension. I.e., we define
\[
G \equiv_{M}^\mathbb{P} H\qquad \iff\qquad M[G]=M[H].
\]

The equivalence relation \(\equiv_{M}^\mathbb{P}\) was first analyzed from the point of view of the Borel complexity theory by Smythe~\cite{Smy22}.
By the definability of the forcing relations, this is a countable Borel equivalence relation on the Polish space $\Gen(M, \mathbb{P})$. (See \cite[Lemma~2.6]{Smy22}.)
The case where \(\mathbb{P}\) is the Cohen forcing \(\mathbb{C}\) is particularly compelling, and it is motivated by a long-standing open problem in descriptive set theory. Recall that a Borel equivalence relation is \markdef{hypefinite} if is the increasing union of a sequence of Borel equivalence relations with finite equivalence classes.
While \cite[Theorem~3.1]{Smy22} proves that \(\equiv_{M}^\mathbb{C}\) is an increasing union of hyperfinite equivalence relations, it is unclear whether \(\equiv_{M}^\mathbb{C}\) is hyperfinite.
This is a particular instance of a central problem in descriptive set theory: the union problem.

The union problem was first posed by Doughety, Jackson, and Kechris~\cite{DouJacKec} in their seminal work about hyperfinite equivalence relation. Some recent striking results show that certain regularity assumption such as Borel asymptotic dimension and \(u\)-amenability  provides a positive solution to the union problem. (See respectively \cite{CJMSTD} and \cite{NarVac}.) However, the union problem remains open and is entangled with another central theme in descriptive set theory: the nuance between hyperfiniteness and measure hyperfiniteness.

On the other hand, Frisch, Shinko and Vidnyanszky~\cite{FriShiVid} recently explored the scenario of the union problem being false and its complexity consequences. 

Back to generic equivalence, people have wondered whether \(\equiv_{M}^\mathbb{C}\) constitutes a counterexample to the union problem. The combinatorial nature of forcing posets suggests that one could develop new techniques to study hyperfiniteness in the more restrictive context of generic equivalence. (For a discussion about this, we refer the reader to~\cite[Page~82]{Kec24}). 

In addition, Smythe~\cite[Question~1]{Smy22} asks to determine the Borel complexity of equivalence of generics for other forcings. In this paper we continue on this theme. We primarily focus on the problem of determining smoothness and hyperfiniteness for \(\equiv^\mathbb{P}_M\).

In Section~\ref{sec : smoothness} we provide a new characterization of smoothness for \(\equiv_{M}^\mathbb{P}\). In Sections~\ref{sec : Prikry}  we prove that \(\equiv_{M}^\mathbb{P}\) is hyperfinite in the case of Prikry. In Section~\ref{sec : Cohen} we prove some additional results about genric equivalence and we discuss some equivalence relations closely related to \(\equiv^\mathbb{C}_M\). We believe that these would be useful to address the problem of finding the Borel complexity of \(\equiv^\mathbb{C}_M\) and other forcing in the future.




\section{Preliminaries}


\subsection{Forcing}





  Recall that $\mathbb{P}$ is \markdef{weakly homogeneous} if for all pairs of conditions $q,r$, there is an automorphism $e\colon\mathbb{P}\rightarrow\mathbb{P}$, such that $e(q)$ is compatible with $r$. Note that this implies that there are generic filters $G,H$, such that $q\in G, r\in H$ and $V[H]=V[G]$. Also, recall that if, on the other hand $G$ and $H$ are two generic filters for $\mathbb{P}$ with  $V[H]=V[G]$, and $\overline{\mathbb{P}}$ is the Boolean completion of the poset, then there is an  automorphism $e\colon\overline{\mathbb{P}}\rightarrow \overline{\mathbb{P}}$, such that $e"G=H$.

Weak homogeneity can be used to detect non-smoothness in the following precise sense. Recall that a poset is \markdef{atomless} if and only if every condition has incompatible extensions. The following result can be found in~\cite[Theorem~2.12]{Smy22}.

\begin{theorem}[Smythe~2022]
\label{thm : Smythe not smooth}
If $\mathbb{P}$ is atomless and weakly homogeneous, then $\equiv_{M}^\mathbb{P}$ is not smooth.
\end{theorem}

In fact, Smythe proved that $\mathbb{P}$ is weakly homogeneous if and only if the action of $\Aut^M(\mathbb{P})$ on $\Gen(M, \mathbb{P})$ is generically ergodic. (See~\cite[Lemma~2.11]{Smy22}.) In Proposition~\ref{prop : weakly homogeneous}  below, we show that the sufficient condition in Theorem~\ref{thm : Smythe not smooth} is not necessary.

We will say that $\mathbb{P}$ is \markdef{densely weakly homogeneous} if for densely many $p$, $\mathbb{P}_{\leq p}$ is weakly homogeneous (i.e., the poset $\{q\in \mathbb{P}\mid q\leq p\}$).
Clearly, weak homogeneity implies densely weak homogeneity.

For an example of densely weakly homogeneous poset that is not weakly homogeneous, take a lottery sum of non isomorphic weakly homogeneous forcings.  Recall that the \markdef{lottery sum} of a collection \(\{\mathbb{P}_i:i\in I\}\) of forcing posets is the
poset \(\bigoplus_{i\in I} P_i = \{\langle  i, p\rangle : i\in I \text{ and } p \in \mathbb{P}_i\}\cup\{\mathbb{1}\}\), ordered with \(\mathbb{1}\) weaker than everything else and
\(\langle i, p\rangle \leq \langle j, q\rangle\) when \(i = j\) and \(p\leq_{\mathbb{P}_i} q\).  
Then for all $\langle i, p\rangle$, the forcing below 
 $\langle i, p\rangle$, the poset $\mathbb{P}\downarrow {\langle i, p\rangle}$ is isomorphic to \(\mathbb{P}_i\downarrow p\), which is just $\mathbb{P}_i$by homogeneity, and therefore weakly homogeneous, and so the full forcing is densely weakly homogeneous. On the other hand, since the $\mathbb{P}_i$'s are no isomorphic, the full forcing is not weakly homogeneous



\subsection{Countable Borel equivalence relations}



An equivalence relation \(E\) is countable if every \(E\)-class is countable. It is finite if every \(E\)-class is finite.

Let \(X = 2^\mathbb{N}\). Then the \markdef{eventual equality relation}
\[
x\mathbin{E}_0y\quad \iff\quad \exists m \forall n \geq m(x_n = y_n).
\]
It is immediate that \(E_0\) is hyperfinite because \(E_0 = \bigcup{i\in\mathbb{N}} F_i\), where \(x\mathbin{F_i} y \iff x_n= y_n\) for all \(n\geq i\). Clearly each \(F_i\) class has cardinality \(2^i\).
\begin{comment}
    
and the \markdef{tail equivalence relation}
\[
x\mathbin{E}_ty \quad \iff\quad \exists m \exists k\forall n(x_{m+n} = y_{k+n})
\]
are countable Borel.

We recall the following result from~\cite{DouJacKec}.

\begin{proposition}[Dougherty-Jackson-Kechris]
 The tail equivalence relation   \(\mathbin{E}_t\) is hyperfinite.
\end{proposition}
\end{comment}




Let \(\Gamma\) be a countable group with the discrete topology.
If \(X\) is a Polish space and \(a\colon \Gamma\times X \to X\) is a continuous action of \(\Gamma\) on \(X\) we say that \(X\) is a \markdef{Polish \(\Gamma\)-space}. In which case, we denote by \(\mathcal{R}(\Gamma\curvearrowright X)\) the corresponding orbit equivalence relation.
Given a Polish \(\Gamma\)-space $X$, a point $x\in X$ is \markdef{condensed} if \(x\) is an accumulation point of its orbit. That is,  for all open neighborhoods $U$ of $x$, $|U\cap[x]|\geq 2$. In the sequel, we make use of the following well-known fact. For a proof see e.g.~\cite[Proposition~2.2]{CalCla23'}.

\begin{fact}
Let \(X\) be a Polish \(\Gamma\)-space.  The orbit equivalence relation \(\mathcal{R}(\Gamma\curvearrowright X)\) is not smooth if and only if there is a condensed point.
\end{fact}





\section{Characterization of smoothness}
\label{sec : smoothness}

The main goal in this section is to present a characterization of non-smoothness for \(\equiv_{M}^\mathbb{P}\). (Theorem~\ref{lemma : dagger} below.)
First we spell out condensation in the case when $X$ is $\Gen(M, \mathbb{P})$ and \(\Gamma=\Aut^M(\mathbb{P})\).
Recall that \(N_p=\{G\in \Gen(M, \mathbb{P})\mid p\in G\}\). It is crucial for us that the set \(\{N_p \mid p\in \mathbb{P}\cap M\}\) is a basis for the space \(\Gen(M,\mathbb{P})\). (For the details see~\cite[Lemma~2.3]{Smy22}.) Therefore we can refer to the basic (cl)open sets \(N_p\), rather than referring to all open subset of the space.

\begin{definition}
\label{def : condensed filter}
An \(M\)-generic filter $G\in \Gen(M, \mathbb{P})$ is \markdef{condensed} if for all $N_p$ with $G\in N_p$, there is $H\neq G$, such that $M[H]=M[G]$ and $H\in N_p$. 
\end{definition}




\begin{definition} Let \(\mathbb{P}\) be a forcing poset in \(M\). For any \(p\in \mathbb{P}\) define the statement \eqref{dwh} as follows:
\begin{multline}\tag{$\dagger_p$}\label{dwh}
\text{For all $p'\leq p$, there are incompatible $q,r\leq p'$,}\\
\text{such that there exists distinct generic filters $G,H$,}\\
\text{such that $q\in G, r\in H$ and $V[H]=V[G]$\}.}
\end{multline}
\end{definition}

In other words condition~\eqref{dwh} says that we can find instances of homogeneity densely often below $p$. Clearly, densely weakly homogeneous implies  $(\dagger)_p$ for all $p$.
We will show that this is strictly weaker than being densely weakly homogeneous.




\begin{theorem}
\label{lemma : dagger}
For any forcing poset \(\mathbb{P}\), we have $\equiv_{M}^\mathbb{P}$ is not smooth if and only if \eqref{dwh} holds for some $p\in \mathbb{P}$.

\end{theorem}
\begin{proof}
For the forward direction,  if $\equiv_{M}^\mathbb{P}$ is not smooth, there is a condensed filter \(G\in\Gen(M,\mathbb{P})\). It follows that \eqref{dwh} holds for all $p\in G$. (Use Definition~\ref{def : condensed filter}.)

For the right to left direction,
assume \eqref{dwh} for some \(p\in\mathbb{P}\).  Let $G$ be a generic filter with $p\in G$.  We want to show that $G$ is condensed.   Namely,  for all \(p'\in G\) there is \(H\neq G\) with  and \(M[G] = M[H]\) and \(H \ni p'\).

Fix \(p'\in G\). By strenghtening if necessary, we may assume $p'\leq p$. First we prove the following density lemma.
\begin{proof}[Claim]
The set \(\mathcal{D} = \{ q : \exists e \in \Aut(\mathbb{P})\cap M\text{ such that }e(q) \perp q, e(q)\leq p' \}\) is dense in \(\mathbb{P}\) below  \(p'\).
\renewcommand{\qed}{}
\end{proof}
\begin{proof}
Given \(q\leq p' \leq p\) by \eqref{dwh} we can find incompatible conditions \(q', r\leq q\) such that
\[
q'\in H_1 \quad\quad r\in H_2 \quad\quad M[H_1] = M[H_2].
\]

Then there is \(e\in \Aut(\bar{\mathbb{P}})\cap M\) such that $e"H_1=H_2$. So, \(e(q')\) and \(r\) are compatible. Then let \(r^* \leq e(q'), r\) and denote by \(q^* = e^{-1}(r^*)\). It follows that \(q^*\leq q'\) and also that \(q^* \) and \(e(q^*)=r^*\) are incompatible. Then $q^*\in \mathcal{D}$.
\end{proof}

Choose \(q\in\mathcal{D}\cap G\) so that \(q\leq p'\).  Let \(e\in \Aut(\mathbb{P})\cap M\) witness that \(q\in\mathcal{D}\). By letting \(H = e'' G\) we get \(M[G] = M[H]\) and since \(q\) and \(e(q)\) are incompatible it follows that \(G,H\) are different. And since both \(q\) and \(e(q)\) are below $p'$, both $G,H$ contain $p'$, as desired.
\end{proof}












Using condensation we show the following:
\begin{proposition}
\label{prop : weakly homogeneous}
There is a poset $\mathbb{P}$ that is not weakly homogeneous,  but \(\equiv_{M}^\mathbb{P}\) is not smooth.
\end{proposition}








Towards the proof of Proposition~\ref{prop : weakly homogeneous}, we will define \(\mathbb{P}_\ell\) such that  \(\mathbb{P}_\ell\) is not densely weakly homogeneous,  but it satisfies \eqref{dwh} for all \(p\in\mathbb{P}_\ell\).
The idea is to take lottery sums of nonisomorphic weakly homogeneous forcings in a tree like fashion. 
 Let $\mathbb{P}_0$ be a weakly homogeneous forcing poset. For ease of exposition we take a constant sequence of forcing \(\{\mathbb{P}^i_0: i\in \mathbb{N}\}\) with \(\mathbb{P}^i_0 = \mathbb{P}_0\). Let \(\{\mathbb{P}_1^i:i\in\mathbb{N}\}\) be a sequence of nonweakly homogeneous forcings such that no \(\mathbb{P}^i_1\) projects\footnote{Here \(\mathbb{P}_0\) projects onto \(\mathbb{P}_1\) means that whenever \(G\) is generic for \(\mathbb{P}_0\), one can define a generic for \(\mathbb{P}_1\) in \(V[G]\).} into \(\mathbb{P}^j_1\) for \(i\neq j\).
For nonempty \(\sigma\in 2^{<\omega}\) let \(\mathbb{P}_\sigma\) be \(\prod_{n<|\sigma|}\mathbb{P}^n_{\sigma(n)}\), so that $\mathbb{P}_{01}=\mathbb{P}_0\times\mathbb{P}^1_1, \mathbb{P}_{11}=\mathbb{P}^0_1\times\mathbb{P}^1_1,\mathbb{P}_{110}=\mathbb{P}^0_1\times\mathbb{P}^1_1\times\mathbb{P}_0$, etc.  Finally we define \(\mathbb{P}_\ell\) as follows:
Conditions in $\mathbb{P}_\ell$ are of the form $(\sigma, p)$ for some $\sigma\in 2^{<\omega}$ and \(p\in \mathbb{P}_\sigma\), and

\[(\sigma, p)\leq (\tau,q)\qquad\iff\qquad\sigma \supseteq \tau\text{ and } {p\restriction \tau} \leq_{\mathbb{P}_\tau} q.\]

\begin{lemma}
\(\mathbb{P}_\ell\) is not densely weakly homogeneous,  but \eqref{dwh} holds for all $p$.
Therefore, \(\equiv_{M,\mathbb{P}_\ell}\) is not smooth.
\end{lemma}
\begin{proof}
Since \(\mathbb{P}^i_1\) are not weakly homogeneous, we can find conditions \(q_1,r_1\in \mathbb{P}^i_1\) such that for all \(\mathbb{P}^i_1\)-generics \(G_1\ni q_1\) and \(H_1\ni r_1\) we have \(V[G_1]\neq V[H_1]\).
It follows that the poset $\mathbb{P}_\ell$ is not weakly homogeneous because any condition $(\sigma,p)$ can be extended to incompatible conditions $({\sigma\cat 1}, p\cat q_1)$ and \(({\sigma\cat 1}, p\cat r_1)
\). Then if \(G,H\) are \(\mathbb{P}_\ell\)-generic containing $({\sigma\cat 1}, p\cat q_1)$ and \(({\sigma\cat 1}, p\cat r_1)\) respectively, and \(V[G]=V[H]\), then  the projections of \(G,H\) on the \(|\sigma|\) coordinate must produce the same model.

On the other hand,  any condition $(\sigma,p)$ can be extended to conditions  $({\sigma\cat 0}, q)$, $({\sigma\cat 0},r)$ such that \(q,r\in\mathbb{P}_{\sigma\cat 0}\) and  $q\upharpoonright |\sigma|=r\upharpoonright |\sigma|=p$.  Using homogeneity of $\mathbb{P}_0$ we can find generics containing these conditions that produce the same extension.  Therefore,  we conclude that \(\mathcal{R}(\Aut(\mathbb{P}_\ell)\curvearrowright \Gen(\mathbb{P}_\ell, M))\) is not smooth by Lemma~\ref{lemma : dagger},  and neither \(\equiv_{M,\mathbb{P}_\ell}\) is.
\end{proof}






\section{Prikry is hyperfinite}
\label{sec : Prikry}

Let $\kappa$ be a measurable cardinal and fix a normal measure $U$ on $\kappa$ (i.e. $U$ is a $\kappa$-complete non principal normal ultrafilter on $\kappa$). In this section $\mathbb{P}$ will be the a Prikry forcing at $\kappa$ with respect to $U$. I.e. conditions are of the form $(s, A)$, where $s$ is a finite increasing sequence of points in $\kappa$ and $A\in U$. We will refer to set in $U$ as ``measure one". For $(s, A)\in\mathbb{P}$, we call $s$ the stem of the condition, and $|s|$ is the length. We say $q$ is a direct extension of $p$, denoted $q\leq^* p$ if $q\leq p$ and they have the same stem.


Let $M$ be a countable model satisfying enough of \textsf{ZFC}, containing $\kappa$ and $\mathbb{P}$.

A generic filter for $\mathbb{P}$ induces a generic Prikry sequence $\vec{\alpha}=(\alpha_n: n<\omega)$ cofinal in $\kappa$, and it is actually defined from it as follows: given such a sequence $\vec{\alpha}$, define $G$ as the set of conditions $( s, A)$, such that for some $n$, $s = ( \alpha_0, \dotsc, \alpha_{n-1})$ and for all $i\geq n$,  $\alpha_i\in A$.  We say that $\vec{\alpha}$ is a \markdef{generic sequence} if the filter induced from it is generic. We will be identifying generic filters with the corresponding generic Prikry sequences.


By identifying $M$ with its transitive collapse, we can assume that \(M\) is transitive without any loss of generality.


\begin{theorem}
\label{thm : Prikry}
$\equiv_{M}^\mathbb{P}$ is hyperfinite.
\end{theorem}

We will use the following fact about Prikry forcing.
\begin{fact}
\label{fact : Prikry}
\begin{enumerate-(1)}
\item
\label{fact : Prikry(1)}
If $\vec{\alpha}$ and $\vec{\beta}$ are both $\mathbb{P}$-generic sequences over $M$, then $M[\vec{\alpha}]=M[\vec{\beta}]$ if and only if $\vec{\alpha}$ and $\vec{\beta}$ coincide on a tail.
\item
\label{fact : Prikry(2)}
(The Prikry lemma). For any condition $p$ and a dense open set $D$, there is a direct extension $q\leq^*p$ and $k<\omega$, such that every $r\leq q$ of length $k$ is in $D$.
\item
Prikry forcing is weakly homogeneous. In particular, for any $\mathbb{P}$-generic $G$ over $M$ and any condition $p\in M$, there is a $\mathbb{P}$-generic $H$ over $M$, such that $p\in H$ and $M[G]=M[H]$.

\end{enumerate-(1)}
\end{fact}



Fact~\ref{fact : Prikry}\ref{fact : Prikry(1)} implicitly follows from the results~\cite{GitKanKoe}.
Since this work is unpublished, for completeness we provide a proof below.

\begin{proof}
One direction is clear. For the other direction, suppose for contradiction that $\vec{\alpha}$ and $\vec{\gamma}$ are two generic Prikry sequences over $M$, such that $M[\vec{\alpha}]=M[\vec{\gamma}]$, but the symmetric difference of these sequences is infinite. It follows that there is a subsequence of $\vec{\gamma}$, which is disjoint from $\vec{\alpha}$. Let $\vec{\beta}$ enumerate this subsequence. Note that $\vec{\beta}\in M[\vec{\alpha}]$.


Let $p$ in the generic filter induced by $\vec{\alpha}$, force that $\dot{\vec{\beta}}$ is disjoint from the generic Prikry sequence. 

For each $n$, let $D_n$ be the dense set of conditions deciding a value for $\dot{\beta}_n$. By the Prikry lemma, there is a direct extension $p'\leq^*p$, and $k<\omega$, such that every $k$-step extension of $p'$ is in $D_n$. I.e., for every $r\leq p'$ with length $\lh(r)=\lh(p)+k$, there is $\delta<\kappa$, such that $r\Vdash \delta=\dot{\beta}_n$.

\begin{claim}
Fix $n$, and suppose that $k$ is minimal such that the above holds for some $p'\leq^*p$ and $D_n$, and denote $p' = (h, A)$. Then there is a measure one set $B\subset A$, such that for all increasing $s$ in $B^k$, if $\delta_s$ is the unique such that $(h\cat s, B\setminus \max(s)+1)\Vdash\delta=\dot{\beta}_n$, we have that $\max(s)<\delta_s$.
\end{claim}
\begin{proof}
Suppose otherwise. Since $p'$ forces that $\dot{\beta}_n$ is not on the Prikry sequence, each $\delta_s\notin s$.  Let $\phi\colon[A]^k\rightarrow \{0,1\}$ be the function defined as follows
\[
\phi(s) = 
\begin{cases}
0 & \text{if $\max(s)<\delta_s$}\\
1 & \text{if $\max(s)>\delta_s$}.
\end{cases}
\]
By measurability, $\phi$ is homogeneous on a measure one set, namely, there is $B\subset A$, $B\in U$, such that for all $k$, $\phi\upharpoonright [B]^k$ is constant. So, for all $k$-length stems $s$ taken from $B$, $\delta_s<\max(s)$. Then for each $t\in [B]^{k-1}$, the function $g_t\colon \alpha\mapsto \delta_{t\cat\alpha}$ is regressive. By Fodor's Theorem (e.g., see~\cite[Theorem~8.7]{Jec}, there is a measure one set $B_t\subset B$ and $\delta_t<\kappa$, such that for all $s$ of the form $t\cat\alpha$ with $\alpha\in B_t$, $\delta_s=\delta_t$. Let $B'=\triangle B_t = \{\alpha\mid \alpha\in\bigcap_{\max(t)<\alpha} B_t\}$. Then if $r\leq (h, B')$ is any condition with length $\lh(p)+k-1$; denote $r = (h\cat t, C)$. We may assume that $\max(t)<\min(C)$.  Then $C\subset B_t$, and so $r\Vdash \dot{\beta}_n=\delta_t$. Contradiction with the minimality of $k$.
\end{proof}



Using the above claim and inductively applying the Prikry lemma\footnote{cf. Fact~\ref{fact : Prikry}\ref{fact : Prikry(2)}} for each $n$, we can find a condition $p' = (h, A)\leq^*p$, and a sequence $( k_n: n<\omega)$, such that:
\begin{enumerate}
\item
for each $n$, every $k_n$-step extension of $p'$ decides the value of $\dot{\beta}_n$,
\item
for each stem $s$ of points in $A$ of length $k_n$, setting $\delta_s$ to be such that $(h\cat s, A\setminus \max(s)+1)\Vdash\dot{\beta}_n=\delta_s$, we have that $\delta_s>\max(s)$.


\end{enumerate}

Now, let $G$ be the generic filter induced by $\vec{\alpha}$.
Since we can do this densely often, we can find such a condition in  $G$. Now go back to $M[G]$. Let $n$, be such that $\beta_n\in A$. Let $q\leq p$ be the (unique) weakest $k_n$-step extension of $p$ in $G$. Denoting $q = (h\cat s, A\setminus \max(s)+1)$, we have that $q\Vdash\dot{\beta}_n=\delta_s$. Since $q\in G$, then $\delta_s=\beta_n\in A$. And since $\max(s)<\delta_s$, we can further extend $q$ to $r = (h\cat {s\cup\{\delta\}}, A\setminus \delta+1)$. But then $r$ forces that $\dot{\beta}_n$ is on the Prikry sequence. Contradiction with our initial assumption.
\end{proof}


\begin{proof}[Proof of Theorem~\ref{thm : Prikry}]
Since \(M\) is countable, we identify \((\kappa\cap M)^\omega\) with the Baire space. Roughly, speaking we encode \(\equiv^\mathbb{G}_M\) into a hyperfinite equivalence relation on the standard Borel space of \(\mathcal{G}\) of generic Prikry sequences in \((\kappa\cap M)^\omega\). Precisely, since
there is a one-to-one Borel function $f\colon\Gen(M,\mathbb{P})\rightarrow (\kappa\cap M)^\omega$ associating to any filter the corresponding generic Prikry sequence, let \(\mathcal{G} = f'' Gen(M,\mathbb{P})\) which is injective image of a Borel set, therefore Borel.

Working in $V$, enumerate $M\cap\kappa$ by $(\gamma_n: n<\omega)$.
Define $E_n$ as follows: for two $M$-generic Prikry sequences $\vec{x}$ and $\vec{y}$ in \(\mathcal{G}\), let $\vec{x}\mathbin{E}_n\vec{y}$ if and only if 
\begin{itemize}
\item
$x_k=y_k$, for all $k\geq n$, and
\item
$\vec{x}=\vec{y}$ or for all $k<n$, $x_k, y_k\in \{\gamma_i\mid i\leq n\}$.

\end{itemize}
Then each $E_n$ is a finite equivalence relation. In fact, \(|[\vec{x}]_{E_n}| \leq n!\). Moreover,  $ E_n\subset E_{n+1}$, and $\equiv_{M}^\mathbb{P} = \bigcup_n E_n$. So $\equiv_{M}^\mathbb{P}$ is hyperfinite, and in particular $\equiv_{M}^\mathbb{P}\sim_B E_0$. 
\end{proof}

\section{Towards hyperfiniteness via mutual genericity}

\label{sec : Cohen}

Let $M$ be the ground model, and let $\mathbb{P}\in M$ be a poset.  Let $X=\Gen(M,\mathbb{P})$ be the set of $M$-generics for $\mathbb{P}$. Note that $X\notin M$, but $X\subset 2^{\mathbb{P}}$ and $2^{\mathbb{P}}\in M$.
Throughout this section, we analyze the equivalence relation $\equiv_{M}^\mathbb{P}$ on $X$ and discuss the interplay between invariant uniformization and hyper-hyperfiniteness. (Invariant uniformization of equivalence relations is also largely discussed in a recent paper by Kechris and Wolman~\cite{KecWol}.) Moreover, we show that a weak form of invariant uniformization implies hyperhyperfiniteness for \(\equiv^\mathbb{P}_M\). (See Proposition~\ref{prop : weak uniformization}.)

We simplify our notation by letting \(E \) be $\equiv_{M}^\mathbb{P}$. 



Following the analysis in Smythe,  we have that $E$ is generated by the countable group $\Aut^M(\overline{\mathbb{P}})$, where $\overline{\mathbb{P}}$ is the Boolean algebra completion of $\mathbb{P}$. (E.g., see \cite[Theorem~2.16]{Smy22}.)
Let $\{\gamma_n\mid n<\omega\}$ enumerate the automorphisms in $\Aut^M(\overline{\mathbb{P}})\).   For each $n$, let $\Gamma_n$ be the group generated by  $\{\gamma_i\mid i\leq n\}$.
Let $E_{\Gamma_n}$ be the induced orbit equivalence relation on $2^{\mathbb{P}}$. (We note that the full \(2^{\overline{\mathbb{P}}}\) may not be a Polish space, so we take the restriction just to $\mathbb{P}$.) Note that while the full enumeration of the automorphisms is not in $M$, for each $n$, both \(\Gamma_n\) and \(E_{\Gamma_n}\) are in \(M\).  Moreover $E_{\Gamma_n}\subseteq E_{\Gamma_{n+1}}$ and $E=\bigcup_n E_{\Gamma_n}$. 




Let $\{I_n\mid n<\omega\}\in M$ be an increasing family of infinite subsets of $\omega$, such that $\bigcup_n I_n=\omega$ and $I_{n+1}\setminus I_n$ is infinite.
Find Borel involutions $\{g_k\mid k<\omega\}$ on $2^{\mathbb{P}}$ (possibly with repetitions), such that for each $n$, $\{g_k\mid k\in I_n\}\in M$ and $E_{\Gamma_n} = \bigcup_{k\in I_n}\text{graph}g_k$. We do this by applying Feldman-Moore for $E_{\Gamma_n}$ in $M$ for each $n$. Then, also $E = \bigcup_k \text{graph}g_k$.


Next we define finite Borel equivalence relations \(E_s\), for \(s\in \omega^{<\omega}\), so that if $s\subseteq s'$, then $E_s\subseteq E_s'$ and each $E_s\subset E$ with a cascade argument as in \cite[Theorem~12.1]{KecMil}.
More precisely, we first fix a Borel linear ordering $<$ on $2^{\mathbb{P}}$. Then, for each $n<\omega$, let $f_k(x) = \min_{<}\{x, g_k(x)\}$. And for $s=\langle s_0, \dotsc, s_{l-1}\rangle\in\omega^{<\omega}$,  set
\[x \mathbin{E_s} z\qquad\iff\qquad f_{s_{l-1}}\circ \dotsb \circ f_{s_1} \circ f_{s_0}(x)=f_{s_{l-1}}\circ \dotsb \circ f_{s_1} \circ f_{s_0}(z).\]
Note that each $E_s\in M$ and the sequence $\langle E_s\mid s\subset I_n\rangle\in M$ for each $n\in \omega$.

Next, fix $n$ and let $z\mathbin{E_{\Gamma_n}} x$ and $s\in I_n^{<\omega}$. Let
\[x' = f_{s_{l-1}}\circ \dotsb \circ f_{s_1} \circ f_{s_0}(x)\qquad\text{ and }\qquad z' = f_{s_{l-1}}\circ \dotsb \circ f_{s_1} \circ f_{s_0}(z)\] where $l=|s|$. Then $x' \mathbin{E_{\Gamma_n}} z'$, and so  $g_k(x')=z'$, for some $k\in I_n$.  Let $s'=s\cat k$. Then we have that $z\mathbin{E_{s'}}x$. Take say the mimimun such $s'$ and denote it by $s^n_{x,z}$. Note that the function  $(x,z,s)\mapsto s^n_{x,z}$ is in \(M\). 



Now, for any $x,z$ such that $x\mathbin{E} z$, and $s\in \omega^{<\omega}$, let $s'_{x,z}=s^n_{x,z}$, where $n$ is the least such that $z\mathbin{E_{\Gamma_n}} x$ and $s\subset I_n$. Of course, the map  $\Phi\colon (x,z,s)\mapsto s'_{x,z}$ is not in $M$, however, $\Phi$ restricted to $s\in I_n^{<\omega}$ and $x,z$ with $z\mathbin{E_{\Gamma_n}} x$ is in $M$. 



For $x,z$, such that $x\mathbin{E}z$, let $D_{x,z}=\{s'_{x,z}\mid s\in \omega^{<\omega}\}$. This is a dense subset of the poset $(\omega^{<\omega}, \supseteq)$ (which can be identified with $\Add(\omega, 1)$). And define: $$C_{x,z}=\bigcup_{s\in \omega^{<\omega}}\mathcal{N}_{s'_{x,z}}.$$
Here we use the notation $\mathcal{N}_s:=\{y\in\omega^\omega\mid s\sqsubset y\}$. Clearly, $C_{x,z}$ is dense open. For each $x$, let $C_x = \bigcap_{z,w\in[x]_E}C_{w,z}$. Then $C_x$ is comeager in $\omega^\omega$. Also, for each $n$ and $x,z$, such that $xE_{\Gamma_n}z$, define $$C^n_{x,z}=\bigcup_{s\in I_n^{<\omega}}\mathcal{N}_{s^n_{x,z}}.$$
Furthermore, for each $x$, let $C^n_x = \bigcap_{z,w\in[x]_{E_{\Gamma_n}}}C^n_{w,z}$. Note that $C^n_x$ is comeager in $I_n^{<\omega}$.



We will say that $y\in\omega^\omega$ is \textbf{\(N\)-generic} for some model \(N\), if it is generic for the poset $(\omega^{<\omega},\supseteq)$ over \(N\).
Recall that $X=\Gen(M,\mathbb{P})$ is the the set of $M$-generics for some poset $\mathbb{P}$.

\begin{definition}
     We say that $f\colon X\to\omega^\omega$ is a {\bf generic  \(\equiv^\mathbb{P}_M\)-invariant uniformization} if and only if
     \begin{enumerate-(i)}
    \item 
    $f(x)$ is $M[x]$-generic for all \(x\in X\); and 
    \item
    \(f\) is a \(\equiv^\mathbb{P}_M\)-invariant uniformization; i.e., \(x\equiv^\mathbb{P}_My\implies f(x)=f(y)\).
     \end{enumerate-(i)}
\end{definition}
 




\begin{theorem}
\label{thm : invariant function}
If there is a generic \(\equiv^\mathbb{P}_M\)-invariant uniformization \(g\colon X\to\omega^\omega\), then the equivalence relation $\equiv ^\mathbb{P}_M$ is hyperhyperfinite.
\end{theorem}
\begin{proof}
Let \(E\) denote \( \equiv ^\mathbb{P}_M\) and  suppose that $g$ is such a function. 


%Define $f\colon \Gen(M, \mathbb{P})\rightarrow \Gen(M, \mathbb{P})\times \omega^\omega$ by $f(x)=(x, g(x))$. Then $f$ witnesses that $E^\mathbb{P}_M$ is Borel reducible to $E^*$ defined by \[(x,y)\mathbin{E^*} (x', y')\qquad\iff\qquad x\mathbin{E}x',\; y=y'\;\text{and } x\mathbin{E_y} x'\].
For each $n$ and $y\in\omega^\omega$, let $E^n_y$ be defined as above, with respect to the Borel involutions generating $E_{\Gamma_n}$.  Now, define $E_n$ by setting \[x\mathbin{E_n}z\qquad\iff\qquad x\mathbin{E}z\quad\text{and}\quad xE^n_{g(x)} z.\] Then $E_n$ is a hyperfinite relation because \(E_n= \bigcup_k F_{n,k}\) where \(x\mathbin{F}_{n,k} z \iff x\mathbin{E} z \text{ and } x\mathbin{E}^n_{g(x)\restriction k}z\).

Moreover, if  $x\mathbin{E_{\Gamma_n}}z$, then since $g(x)$ is $M[x]=M[z]$-generic and $D^n_{x,z}\in M[x]$, we have that $x\mathbin{E^n_{g(x)}}z $. It follows that $E^\mathbb{P}_M=\bigcup_n E_n$, and for each $n$, $E_n\subseteq E_{n+1}$. 
\end{proof}

\begin{corollary}
\label{cor : Random}
    For any model \(M\), there is no generic \(\equiv^\mathbb{B}_M\)-invariant uniformization, when \(\mathbb{B}\) is the random forcing. 
\end{corollary}
\begin{proof}
The equivalence relation \(\equiv^\mathbb{B}_M\) is not Fr\'echet amenable by \cite[Theorem~4.4]{Smy22}.
Since every hyperfinite equivalence relation is amenable, every hyper-hyperfinite equivalence relation is Fr\'echet amenable (e.g., see~Proposition~9.11(v) of~\cite{Kec24}).
    Therefore the corollary is an immediate consequence of Theorem~\ref{thm : invariant function}.
    \end{proof}

It turns out that a somewhat weaker assumption is enough to get the conclusion of  Theorem~\ref{thm : invariant function}. Denote \(E_0(\mathbb{N})\) be the equivalence relation of eventaul equality on \(\omega^\omega\). It is well known that \(E_0(\mathbb{N})\sim_BE_0\), therefore \(E_0(\mathbb{N})\) is hyperfinite. (E.g., see Gao~\cite[Proposition~6.1.2]{Gao}.)

\begin{proposition}
\label{prop : weak uniformization}
Suppose that $f\colon X\to \omega^\omega$ is a Borel function such that,
\begin{enumerate}
\item
$f(x)$ is \(M[x]\)-generic, for all \(x\in X\);
\item
$x\equiv^\mathbb{P}_M y$ implies $({f(x)},{f(y)})\in E_0(\mathbb{N})$.

\end{enumerate}
Then $\equiv^\mathbb{P}_M$ is hyperhyperfinite.

\end{proposition}
\begin{proof}
Let $f$ be such a function. As before, construct $E_{\Gamma_n}$, for $n<\omega$,  such that each $E_{\Gamma_n}\in M$, and $E$ is their increasing union. Fix a partition $\{I_n\mid n<\omega\}$ in $M$ of $\omega$ into infinite sets.

For each $n$, let $\{g^n_i\mid i<\omega\}\in M$ be Borel involutions, so that $E_{\Gamma_n}$ is the union of their graph; and for $s\in\omega^{<\omega}$, we define $x\mathbin{E^{\Gamma_n}_s}z$ iff $g_s(x)=g_s(z)$, where $g_s$ is defined as in the cascade argument from as in \cite[Theorem~12.1]{KecMil} mentioned above. Similarly, define $$D^{\Gamma_n}_{x,z} = \{s\mid x\mathbin{E^{\Gamma_n}_s}z\}.$$ We have that this set is dense open for all $x,z$ such that $x\mathbin{E_{\Gamma_n}}z$.

For $n,k<\omega$, define $xF^n_kz$ iff
 \begin{itemize}
 \item
 for all $n'\geq n$, $f(x)\upharpoonright I_{n'} = f(z)\upharpoonright I_{n'}$, and
 \item
 setting $\bar{y} = f(z)\upharpoonright I_n$, we have $x\mathbin{E_{\bar{y}\upharpoonright k}^{\Gamma_n}}z$.
\end{itemize}
Then $\{ F^n_k\mid k<\omega\}$ is an increasing sequence of finite equivalence relations.

Next we check that $\bigcup_{k,n} F^n_k=E$. Suppose that
$xEz$. Then for some $n\in\omega$, $x\mathbin{E_{\Gamma_n}}z$, and $D^{\Gamma_n}_{x,z}\in M[x]=M[z]$. By increasing $n$ if necessary, we may assume that for all  $n'\geq n$, $f(x)\upharpoonright I_{n'} = f(z)\upharpoonright I_{n'}$. Let  $\bar{y} = f(x)\upharpoonright I_n$. Since $f(x)$ is $M[x]$-generic and $I_n$ is infinite, by properties of the Cohen forcing, we have that  $\bar{y}$ is $M[x]$-generic. So for some $k$, $\bar{y}\upharpoonright k\in D^{\Gamma_n}_{x,z}$. It follows that $xF^n_kz$.

Finally, we  check that $\bigcup_k F^n_k=\bigcup_k F^{n+1}_k$:
suppose that for some $k$, $xF^n_kz$. It follows that for all $n'\geq n$, $f(x)\upharpoonright I_{n'} = f(z)\upharpoonright I_{n'}$ and that $x\mathbin{E_{\Gamma_n}}z$. Hence $x\mathbin{E_{\Gamma_{n+1}}}z$. By a similar argument as above we get that for some $k'$, $x\mathbin{F^{n+1}_{k'}}z$.
\end{proof}

\begin{remark}
Edward Hou recently showed that a partial converse to the above proposition is also true \cite{Hou}. Namely, he proved that if $E$ is hyperfinite, then one can find a function as in the assumption of the proposition.
\end{remark}

\begin{proposition}
Suppose that $y\in \omega^\omega$ and $x\in X$ are $M$-mutually generic. (I.e. $y\times x$ is $M$-generic for $\omega^{<\omega}\times\mathbb{P}$.) Suppose also that $\{g_n\mid n<\omega\}\in M[y]$ are Borel involutions, such that $\equiv^\mathbb{P}_M=\bigcup_n \graph(g_n)$. Then $[x]_{\equiv^\mathbb{P}_M} = [x]_{E_y}$. Namely,  $\equiv^\mathbb{P}_M$ restricted to $\Gen(\mathbb{P}, M[y])$ is hyperfinite. 
\end{proposition}
\begin{proof}
As above, denote $\equiv^\mathbb{P}_M$ by \(E\). The inclusion is $[x]_E \supseteq [x]_{E_y}$ is clear. For the nontrivial direction, suppose $z\in X$ is such that $z\mathbin{E}x$. Then $y$ is $M[z]$ generic (as $M[x]=M[z]$), $D_{x,z}\in M[z]$, and so for some $s\in D_{x,z}$, $s\sqsubset y$. It follows that $z\in[x]_{E_s}\subseteq [x]_{E_y}$. 
The hyperfiniteness of $E$ restricted to $\text{Gen}(\mathbb{P}, M[y])$ follows since every  $\text{Gen}(\mathbb{P}, M[y])$ is mutually generic with $y$.
\end{proof}  


Next we look at the Cohen forcing $\mathbb{C}=\Add(\omega, 1)$ to add one real. %We note that in \cite[Theorem~3.1]{Smy22}, it was shown that  $\equiv_{M}^\mathbb{C}$ is hyper-hyperfinite, i.e. an increasing union of hyperfinite equivalence relations.
Recall that this poset is isomorphic to  $(\omega^{<\omega}, \supseteq)$. Let $h\colon 2^{<\omega}\rightarrow\omega^{<\omega}$ code that isomorphism, $h\in M$. Note that then for any $M'\supseteq M$, $x\in 2^\omega$ is $M'$-generic iff $h"x$ is $M'$-generic. In slight abuse of notation we use $h(x)$ to denote $h"x$.

Let \(\mathrm{evens}, \mathrm{odds}\) denote the sets of even and odd natural numbers respectively. For $x\in 2^\omega$, let  $x_l = x\upharpoonright \mathrm{evens}$
and  $x_r = x\upharpoonright \mathrm{odds}$.
Note that if $x$ is $M$-generic, then both $x_l$ and $x_r$ are also $M$-generic and moreover, they are mutually generic, i.e., $x_l$ is $M[x_r]$-generic and vice versa.
Finally, Let $h\colon2^{\mathrm{evens}}\to\omega^\omega$ be some isomorphism in $M$. Note that if $x\in 2^{\mathrm{evens}}$ is $M$-generic, and $y=h(x)$, then so is $y$.

Let \(E\) denote \(\equiv_M^{\mathbb{C}}\), By Feldman-Moore, $E$ is induced by a countable group generated by some Borel involutions $\{g_n\mid n<\omega\}$. Let $N\supset M$ be a countable model, such that the
$\{g_n\mid n<\omega\}\in N$.

Define the following  subequivalence relations of $E$ (both defined on $X$):
\begin{itemizenew}
\item 
$x\mathbin{E^*}z$ iff $x_l=z_l$ and $x_r E z_r$. 
\item
$x E^N z$ iff $x=z$ or
$x Ez$,  $x_l=z_l$, and $x$ is $N$ -generic.
\end{itemizenew}
Note that if there is a generic \(\equiv_M^\mathbb{P}\)-invariant uniformization, then $E\leq_B E^*$. 


\begin{lemma}
$E^{\text{N}}$ is hyperfinite.
\end{lemma}
\begin{proof}
We first make some observations. First note that for every $x\in X$, $x_l$ and $x_r$ are mutually generic. Suppose now that $xE^N z$, $x\neq z$, and denote $w=x_l=z_l$. Then $M[x] = M[w][x_r]=M[w][z_r]=M[z]$. So $x_r \equiv_{M[w]}^{\mathbb{P}}z_r$. By assumption $x$ and therefore $z$ are also $N$-generic. It follows that $x_r, z_r$ are $N$-generic, and $w$ is $N[x_r] = N[z_r]$-generic.


Define a sequence of equivalence relations $E_n, n<\omega$ as follows. Set 
\[x\mathbin{E}_n z \qquad\iff\qquad x\mathbin{E}^Nz \text{ and }x_r \mathbin{E}_{y\upharpoonright n} z_r,\]
for \(y=h(x_l)\).
Then the \(E_n\)'s are increasing finite equivalence relations. To show that their union is $E^N$, suppose $x\mathbin{E}^Nz$. Let   $w=x_l=z_l$ and $y=h(w)$. Then the dense set $D_{x_r, z_r}$ is in $N[x_r]$ and by mutual genericity, $y$ is $N[x_r]$-generic. It follows that  $y\upharpoonright n\in D_{x_r, z_r}$ for some $n$, and so $x_r \mathbin{E}_{y\upharpoonright n} z_r$.
\end{proof}


We end this section with some questions, that may be stepping stones for settling whether equivalence of generics is hyperfinite for the Cohen forcing and other posets.

\begin{question}
    For which posets $\mathbb{P}$, is there a generic \(\equiv^\mathbb{P}_M\)-invariant uniformization? Is there such a uniformization for the Cohen poset?
\end{question}

Note that by Corollary~\ref{cor : Random} there is no invariant mutually generic function for the Random forcing.  On the other hand, if  $\equiv_M^{\mathbb{P}}$ is smooth, then one can define such a function:

\begin{proposition}
     Suppose that  $\equiv_M^{\mathbb{P}}$ is smooth. Then there is a  generic \(\equiv^\mathbb{P}_M\)-invariant uniformization.
\end{proposition}
\begin{proof}
    Let $f\colon\Gen(\mathbb{P}, M) \rightarrow \Gen(\mathbb{P}, M)$ be a Borel selector, and let $\langle\dot{D}_n\mid n<\omega\rangle$ enumerate all $M$-names for dense sets, and finally let $\prec$ be some wellorder of $\mathbb{P}$. Define $g\colon \Gen(\mathbb{P}, M) \rightarrow \Gen(\mathbb{P}, M)$ as follows. For each $x\in \Gen(\mathbb{P}, M)$, inductively construct $\{p_n\mid n<\omega\}$, so that each $p_n$ is the $\prec$-least condition such that $p_n\leq p_{n-1}$ and $p_n\in \dot{D}_n[f(x)]$. Then set $g(x) = \{p\mid \exists n\, (p_n\leq p)\}$. Then $g$ is the desired generic \(\equiv^\mathbb{P}_M\)-invariant uniformization.
\end{proof}


\begin{question}
    Suppose that $\equiv_M^{\mathbb{P}}$ is hyperfinite. Must there be generic $\equiv^\mathbb{P}_M$-invariant uniformization?
\end{question}





\newcommand{\etalchar}[1]{$^{#1}$}
\begin{thebibliography}{CJM{\etalchar{+}}23}

\bibitem[CC24]{CalCla23'}
Filippo Calderoni and Adam Clay.
\newblock Condensation and left-orderable groups.
\newblock {\em Proc. Amer. Math. Soc. Ser. B}, 11:579 -- 588, 2024.

\bibitem[CJM{\etalchar{+}}23]{CJMSTD}
Clinton~T. Conley, Steve~C. Jackson, Andrew~S. Marks, Brandon~M. Seward, and Robin~D. Tucker-Drob.
\newblock {Borel asymptotic dimension and hyperfinite equivalence relations}.
\newblock {\em Duke Mathematical Journal}, 172(16):3175 -- 3226, 2023.

\bibitem[DJK94]{DouJacKec}
Randall Dougherty, Stephen Jackson, and Alexander~S. Kechris.
\newblock The structure of hyperfinite borel equivalence relations.
\newblock {\em Transactions of the American Mathematical Society}, 341(1):193--225, 1994.

\bibitem[FSV24]{FriShiVid}
Joshua Frisch, Forte Shinko, and Zoltan Vidnyanszky.
\newblock Hyper-hyperfiniteness and complexity.
\newblock Preprint available on the arXiv at \url{https://arxiv.org/abs/2409.16445}, 2024.

\bibitem[Gao09]{Gao}
Su~Gao.
\newblock {\em Invariant descriptive set theory}, volume 293 of {\em Pure and Applied Mathematics (Boca Raton)}.
\newblock CRC Press, Boca Raton, FL, 2009.

\bibitem[GKK20]{GitKanKoe}
Moti Gitik, Vladimir Kanovei, and Peter Koepke.
\newblock Inter- mediate models of prikry extensions.
\newblock Preprint, available at \url{https://www.researchgate.net/publication/228378916_Intermediate_models_of_Prikry_generic_extensions}, 2020.

\bibitem[Hou]{Hou}
Edward Hou.
\newblock Personal communication.

\bibitem[Jec03]{Jec}
Thomas Jech.
\newblock {\em Set theory}.
\newblock Springer Monographs in Mathematics. Springer-Verlag, Berlin, 2003.
\newblock The third millennium edition, revised and expanded.

\bibitem[Kec24]{Kec24}
Alexander~S. Kechris.
\newblock {\em The Theory of Countable Borel Equivalence Relations}.
\newblock Cambridge Tracts in Mathematics. Cambridge University Press, 2024.

\bibitem[KM04]{KecMil}
Alexander~S. Kechris and Benjamin Miller.
\newblock {\em Topics in orbit equivalence relations}.
\newblock Lecture notes in Mathematics. Springer-Verlag Berlin Heidelberg, 2004.

\bibitem[KW25]{KecWol}
Alexander~S. Kechris and Michael Wolman.
\newblock Invariant uniformization.
\newblock Preprint available on the arXiv at \url{https://arxiv.org/abs/2405.15111}, 2025.

\bibitem[NV25]{NarVac}
Petr Naryshkin and Andrea Vaccaro.
\newblock Hyper-u-amenablity and hyperfiniteness of treeable equivalence relations.
\newblock Preprint available on the arXiv at \url{https://arxiv.org/abs/2507.07891}, 2025.

\bibitem[Smy22]{Smy22}
Iian~B. Smythe.
\newblock Equivalence of generics.
\newblock {\em Arch. Math. Log.}, 61(5–6):795–812, jul 2022.

\end{thebibliography}



\end{document}
